\section{Structural Emergence of Mean Bias}

The section shows coherent mean bias observed in Transformer models is not an accidental artifact but a predictable consequence of the interaction between language statistics and the specific geometry of high-dimensional operators. We argue that this bias emerges through a three-stage causal chain: frequency-weighted initialization, architectural regeneration, and residual accumulation.

\paragraph{Statistical Origin: Frequency-Weighted Embedding Expectation.}

A ``seed'' bias is introduced at the very first layer of the model.
Let $\mathcal{V}$ denote the vocabulary and let $p(v)$ be the token frequency
distribution over $\mathcal{V}$.
In natural language corpora, $p(v)$ is highly non-uniform and typically
follows a Zipfian distribution.
The expected embedding vector over the corpus is

\begin{equation}
\mu_{\text{embed}}
=
\mathbb{E}_{v \sim p}[E_v]
=
\sum_{v \in \mathcal{V}} p(v) E_v,
\label{eq:embedding_mean}
\end{equation}

where $E_v \in \mathbb{R}^H$ denotes the embedding vector of token $v$
\cite{mikolov2013efficient,mikolov2013distributed,arora2017simple}.

Because high-frequency tokens (e.g., articles, punctuation, and syntactic
markers) dominate $p(v)$ and receive disproportionately many gradient updates
during training, their embedding vectors tend to align along common predictive
directions.
In high-dimensional spaces, such coherent alignment induces a dominant
rank-one component in the embedding matrix even before subsequent Transformer
processing begins.

\paragraph{The Mechanism of Mean Regeneration.}

Even if internal activations are perfectly centered at zero,
the fundamental operators of the Transformer architecture
systematically regenerate a coherent mean.
This regeneration is not incidental but arises from
the algebraic structure of nonlinearities,
attention mechanisms, and residual integration.

\noindent \emph{Non-Odd Nonlinearities as Bias Generators:} Most modern activation functions $\phi(\cdot)$ used in
Transformer architectures—such as ReLU, GELU, and SwiGLU—
are non-odd functions.
For any symmetric zero-mean random variable $z$ satisfying
$\mathbb{E}[z] = 0$ and $z \overset{d}{=} -z$,
non-odd nonlinearities satisfy

\begin{equation}
\mathbb{E}[\phi(z)] > 0,
\label{eq:non_odd_bias}
\end{equation}

under mild regularity conditions
\cite{timkey2021rogue,puccetti2022frequency}.

For example, ReLU is defined as $\phi(z) = \max(0,z)$.
If $z$ is symmetric, then

\begin{equation}
\mathbb{E}[\max(0,z)]
=
\int_{0}^{\infty} z\, p(z)\, dz
>
0,
\end{equation}

because all negative mass is truncated.
Thus ReLU introduces a strictly positive drift even when
the input is perfectly centered.
The same phenomenon holds approximately for GELU and SwiGLU,
which suppress negative regions asymmetrically.
Only strictly odd nonlinearities such as $\tanh$
preserve zero mean in expectation.

Therefore, unless $\phi$ is explicitly odd,
mean regeneration is mathematically unavoidable.

\noindent \emph{Softmax Attention as Exponential Mean Amplification:}
Softmax reinforces mean drift by exponentially suppressing cross-cluster
mass in the attention weights, which contracts variance while preserving
the mixture mean. A compact statement is:

\begin{theorem}[Exponential Mean-Bias Amplification]
\label{thm:exponential_amplification}
Assume a cluster decomposition $x_j=\mu_{c(j)}+\varepsilon_j$ with mixture mean
$\mu=\sum_m w_m\mu_m\neq 0$, and a logit gap
$\langle x_i,x_j\rangle-\langle x_i,x_{j'}\rangle\ge\Delta>0$ for
$i,j\in C_m$, $j'\notin C_m$. Let
\[
a_{ij}=\frac{\exp(\langle x_i,x_j\rangle/\tau)}{\sum_k \exp(\langle x_i,x_k\rangle/\tau)},
\qquad y_i=\sum_j a_{ij}x_j.
\]
Define $R_X^2=\|\bar x\|^2/(\frac1N\sum_i\|x_i\|^2)$ and
$R_Y^2=\|\bar y\|^2/(\frac1N\sum_i\|y_i\|^2)$. Then
\[
R_Y^2-R_X^2 \ge
c\,\frac{\|\mu\|^2\sigma^2}{(\|\mu\|^2+\sigma^2)^2}\,(1-e^{-\Delta/\tau}),
\]
for a constant $c>0$ depending only on cluster balance, where
$\sigma^2=\frac1N\sum_i\|\varepsilon_i\|^2$.
Thus softmax attention strictly amplifies mean bias whenever $\Delta>0$.
\end{theorem}

\paragraph{Residual Integration and High-Dimensional Magnification.}

Mean drift generated by nonlinearities and attention
is integrated through residual connections:

\begin{equation}
\mu_{l+1}
=
\mu_l + \Delta\mu_l,
\label{eq:residual_accumulation}
\end{equation}

where $\Delta\mu_l$ denotes the mean generated by layer $l$.

This additive coupling causes bias to accumulate approximately linearly
with depth when $\Delta\mu_l$ maintains consistent orientation.

Crucially, high dimensionality magnifies this effect.
Let $\bar{\mu}$ denote the average coordinate-wise bias.
Then the Euclidean norm of the mean vector scales as

\begin{equation}
\|\mu\|_2
\sim
\sqrt{H}\,\bar{\mu},
\label{eq:mean_scaling}
\end{equation}

where $H$ is the hidden dimension
\cite{tropp2012user,halko2011randomized}.

In a space with $H=4096$,
a seemingly small per-coordinate bias is amplified
by a factor of $\sqrt{4096}=64$ in total magnitude.
This geometric magnification explains why mean-dominated components
account for a large fraction (empirically $85\%$–$90\%$)
of extreme activation outlier energy.

Thus, mean regeneration arises from:
(i) non-odd nonlinearities,
(ii) convex aggregation in softmax attention,
and (iii) residual accumulation amplified by high dimensionality.


\paragraph{Empirical Verification.}

We verify the above claims using intermediate activation matrices
from a Qwen-0.6B model, collected at early (10k steps) and late (170k steps)
training stages across multiple layers.
We measure mean magnitude using the normalized ratio

\begin{equation}
R = \frac{\|\mu\|_2}{\sqrt{\mathbb{E}\|X\|^2/N}}.
\end{equation}

\noindent \textbf{(1) Mean-bias energy rises over training and deepens with depth.}
We decompose activations into mean + spike + tail components and track the
mean energy fraction across embedding, shallow, middle, and deep layers.
Figure~\ref{fig:energy_grid} shows a consistent increase in mean-dominated energy
from early to late training, with shallow layers exhibiting strong mean bias
early and deep layers becoming pronounced later.
Across all layers, the mean-bias energy fraction grows as training progresses.

\noindent \textbf{(2) Embedding-layer dominance at late training.}
From the embedding column in Figure~\ref{fig:energy_grid}, the embedding outputs
already exhibit a large mean-bias energy fraction, especially at the late stage.
This aligns with the frequency-weighted origin discussed earlier.

\noindent \textbf{(3) Nonlinear operators amplify mean bias.}
We isolate operator effects by visualizing distributions before and after
nonlinear blocks. In early training, the attention softmax intensifies mean bias,
and in late training, the FFN SwiGLU further amplifies it. The operator outputs
consistently show higher mean-energy dominance than their inputs.

\begin{figure}[t]
\centering
\begin{subfigure}{0.49\textwidth}
\includegraphics[width=\textwidth]{figures/energy/layer0_attn_input_output_energy_10k_notitle.png}
\caption{Softmax attention (10k)}
\end{subfigure}\hfill
\begin{subfigure}{0.49\textwidth}
\includegraphics[width=\textwidth]{figures/energy/layer0_ffn_input_output_energy_combo.png}
\caption{FFN SwiGLU}
\end{subfigure}
\caption{Operator-level input/output energy comparison. Left: attention softmax at 10k. Right: FFN SwiGLU. Both operators increase mean-bias energy from input to output.}
\label{fig:operator_change}
\end{figure}

\noindent \textbf{(4) Residual connections accumulate mean bias.}
Consistent with the residual accumulation mechanism described in
Section~2, mean-biased components propagate and add across layers.
Residual paths preserve the coherent component while integrating
newly generated bias, leading to progressively higher mean-energy
fractions as depth and training advance.
